%% Dans ce fichier, les �ventuels accents sont cod�s
%% en ISO-8859-15 (aka latin9)
%% La ligne suivante donne trois exemples 
%% e accent aigu : �, e accent grave : �, e-dans-l'o : �
%%%%%ENTREES%%%%
\begin{tabarg}
\multicolumn{3}{l}{\titretab {Entr\'ees}}\\
\hline
$U$ & Vecteur, \texttt{float} & 1�re composante du vent (unit� quelconque)\\
$V$ & Vecteur, \texttt{float} & 2�me composante du vent (m\^eme unit� que $U$)\\
$F_e$ & Coeff., \texttt{float} & Fr�quence d'�chantillonnage (Hz)\\
$T_{int}$ & Coeff., \texttt{float} & P�riode d'int�gration pour le calcul de la moyenne (s)\\
\end{tabarg}

\medskip
%%%%%SORTIES%%%%
\begin{tabarg}
\multicolumn{3}{l}{\titretab {Sorties}}\\
\hline
$U_r$ & Vecteur, \texttt{float} & Composante du vent dans l'axe du vent moyen (m\^eme unit� que $U$ et $V$)\\
$V_r$ & Vecteur, \texttt{float} & Composante du vent transverse (� gauche) � l'axe du vent moyen (m\^eme unit� que $U$ et $V$)\\
\end{tabarg}
